% Fixture for evals/02: an ERC-style plan moved to the BSF without adapting it.
% Format and required items are compliant, so the checker is quiet; the
% failures are judgment-level. Every name and result is invented.
\documentclass[12pt,a4paper]{article}
\usepackage[margin=2cm]{geometry}
\linespread{1.5}
\begin{document}
\section*{Research Plan}

\section{Subject and scientific background}
Every river-quality model assumes that a sediment community returns to the same state after
rewetting. That assumption fails in the two datasets that track the same sites across several
droughts, where recovery after the third drought takes twice as long as after the first. No
current assay separates community composition from activity at the timescale of rewetting,
which is why the discrepancy has never been tested.

\section{Objectives and significance}
Our vision is a theory of ecological memory that unifies microbial, plant and animal
communities and transforms how ecosystems are managed worldwide.
O1. Determine whether recovery time depends on drought history.
O2. Identify the taxa whose activity carries that dependence.
O3. Establish a universal law of community memory across all ecosystems.

\section{Methodology and plan of operation}
The Israeli group at the Weizmann Institute runs the field campaign and the activity assays.
The U.S. group at Michigan builds the population models. We will apply a new single-cell
activity assay to cores from four sites, fit recovery models, and test O1 and O2.
The Israeli group and the U.S. group will hold joint meetings twice a year.

\section{Risk analysis and alternative paths}
The approach is ambitious and its outcome is uncertain.

\section{Available resources}
Both groups have the necessary equipment and personnel.

We used generative AI only to check grammar.
\end{document}
